\section*{Chapter \ref{ch:mx:decomposing-the-spread} --- Decomposing the Spread}

\begin{solution}{exo:mx:decomposing-the-spread:1}
The mid is 100.02. Quoted half-spread 0.02; effective half-spread $100.03-100.02=0.01$; price improvement $100.04-100.03=0.01$ per share; realised
half-spread $100.03-100.025=0.005$; impact $100.025-100.02=0.005$. The effective spread is $2\times0.01/100.02$, 2.0 basis points.
\end{solution}

\begin{solution}{exo:mx:decomposing-the-spread:2}
With $\varepsilon=-1$: effective $-(49.98-50.00)=0.02$, realised $-(49.98-49.97)=-0.01$, impact $-(49.97-50.00)=0.03$. The provider bought at 49.98
something worth 49.97 five minutes later: she lost a cent per share on the trade, before fees and rebates.
\end{solution}

\begin{solution}{exo:mx:decomposing-the-spread:3}
The quotes move by $\lambda S/2=0.6\times0.02=0.012$ in the trade's direction; $(1-\lambda)S/2=0.008$ is left for processing (in the basic form, the
0.012 is adverse selection and inventory together).
\end{solution}

\begin{solution}{exo:mx:decomposing-the-spread:4}
Implied spread $2(\phi+\theta)=0.04$; information share $0.012/0.020=0.6$. A buy after a buy: $(\phi+\theta)-(\phi+\rho\theta)=\theta(1-\rho)=0.0084$. A buy
after a sale: $(\phi+\theta)+(\phi+\rho\theta)=2\phi+\theta(1+\rho)=0.0316$. The first buy was partly expected (signs are autocorrelated), so it carries less
news, and the bounce from bid to ask adds $2\phi$ to the second.
\end{solution}

\begin{solution}{exo:mx:decomposing-the-spread:5}
No: the difference, 0.28 ticks, is within two standard errors of the five-minute value. The standard errors grow by 2.28 from one to five minutes, close
to $\sqrt5=2.24$, as the proposition predicts when the efficient price's own moves dominate. Once they are removed, the impact is 0.46 ticks at one
minute and still 0.46 at five (standard errors 0.04): the information was absorbed within the minute and stayed.
\end{solution}

\begin{solution}{exo:mx:decomposing-the-spread:6}
In logs, $\beta=\ln(100.10/99.94)^2+\ln(100.12/99.96)^2=5.12\times10^{-6}$ and $\gamma=\ln(100.12/99.94)^2=3.24\times10^{-6}$. Then $\alpha=0.00112$ and
$S=2(e^\alpha-1)/(1+e^\alpha)\approx\alpha$: 11.2 basis points.
\end{solution}

\begin{solution}{exo:mx:decomposing-the-spread:7}
Informed trades: effective 0.54 ticks, impact 2.11, realised $-1.57$. Noise trades: effective 0.53, impact 0.045, realised 0.49. The informed are 21\% of
the volume: $0.21\times(-1.57)+0.79\times0.49\approx0.06$, the overall 0.056. \omterm{def:mx:why-there-is-a-spread:informed}{Noise traders} pay the margin and the \omterm{def:mx:why-there-is-a-spread:informed}{informed traders}' profit.
\end{solution}

\begin{solution}{exo:mx:decomposing-the-spread:8}
Over five minutes the mid moves with everything that trades after the fill, including the algorithm's own later child orders pushing the same way: that
is its impact, not information (the metaorder's trades overlap). The five-minute average is also noisy (exercise~5). Measure at shorter horizons,
exclude the algorithm's own later trades or compare with matched flow from other clients before concluding that the flow is informed.
\end{solution}

\begin{solution}{pb:mx:decomposing-the-spread:1}
\textbf{1.} $\varepsilon(p-m_t)=\varepsilon(p-m_{t+h})+\varepsilon(m_{t+h}-m_t)$: add and subtract $\varepsilon m_{t+h}$.
\textbf{2.} 1.04 ticks quoted; 0.535 ticks effective half-spread.
\textbf{3.} $\E[\varepsilon_t(P^\ast_t-m_t)]$, the provider's expected loss to what the aggressor knew: 0.449 ticks, 84\%.
\textbf{4.} 2.05 ticks for informed trades, 0.025 for noise.
\textbf{5.} 0.03, 0.08, 0.35 and 0.48 ticks.
\textbf{6.} See the proposition: split the impact at $P^\ast_{t+h}$ and $P^\ast_t$; the martingale term has mean zero and variance $\sigma^2h$.
\textbf{7.} Noise: the standard error grows from 0.07 to 0.15 ticks, as $\sqrt h$.
\textbf{8.} 0.46 ticks from one minute on (0.85 of the effective half-spread): it subtracts $\varepsilon(P^\ast_{t+h}-P^\ast_t)$, which requires the efficient
price.
\textbf{9.} 0.50 and 0.056 ticks: a tenth of what she charges.
\textbf{10.} One minute.
\textbf{11.} Difference $p_t=\mu_t+\phi x_t+\xi_t$ and substitute $\mu_t-\mu_{t-1}=\theta(x_t-\rho x_{t-1})+u_t$.
\textbf{12.} Huang--Stoll $\lambda=0.09$ (spread 1.07 ticks), MRR 0.13 ($\rho=0.33$), Glosten--Harris 0.10 for 100 shares.
\textbf{13.} They assume the quotes absorb a trade's information by the next trade; here the mid moves toward $P^\ast$ through quote updates over tens of
seconds.
\textbf{14.} 0.14, 0.28, 0.49, 0.56 and 0.48: more lags cover more of the adjustment, then add estimation noise.
\textbf{15.} The \omterm{def:mx:why-there-is-a-spread:informed}{informed traders}' rule depends on the size of the gap $P^\ast-m$, a nonlinearity a regression on signs averages away.
\textbf{16.} The range holds the spread once and volatility that grows with the window; comparing one-day and two-day ranges separates them.
\textbf{17.} Corwin--Schultz 0.6 and 2.0 basis points, Abdi--Ranaldo 0.0 and 3.3, against 1.07.
\textbf{18.} When the spread is a visible share of the bar's range (less liquid instruments, wider spreads), and after checking on a period with known spreads.
\textbf{19.} \emph{Named result}: against a true share of 0.84, the realised-spread method at one minute gives 0.89 (0.85 net of the efficient price's move),
Glosten--Harris 0.10, Huang--Stoll 0.09, MRR 0.13 and the VAR 0.14 to 0.56 depending on the lags; the impact settles at about one minute, and beyond it
the efficient price's own moves add noise that grows as $\sqrt h$.
\textbf{20.} Mostly the \omterm{def:mx:why-there-is-a-spread:informed}{noise traders}, who paid for the \omterm{def:mx:why-there-is-a-spread:informed}{informed traders}' 1.57 ticks per share and left the provider 0.056.
\end{solution}

\begin{solution}{iq:mx:decomposing-the-spread:1}
Effective: $\varepsilon(p-m_t)$, what the trade paid over the mid before it. Realised: $\varepsilon(p-m_{t+h})$, what the provider kept at the later mid. The
difference is the trade's price impact at $h$, mostly adverse selection when $h$ is well chosen.
\iqlookfor{The sign convention; the identity; a horizon.}
\end{solution}

\begin{solution}{iq:mx:decomposing-the-spread:2}
Long enough for quotes to absorb the trade's information, short enough that the price's own volatility and later trades do not swamp it. Too short:
information is counted as spread earned. Too long: noise (standard error growing as $\sqrt h$) and contamination by later trades, including the same
metaorder. Plot the impact against the horizon with standard errors and pick where it flattens.
\iqlookfor{Both failure modes; a data-driven choice.}
\end{solution}

\begin{solution}{iq:mx:decomposing-the-spread:3}
$p_t=\mu_t+\phi x_t+\xi_t$, $\mu_t=\mu_{t-1}+\theta(x_t-\rho x_{t-1})+u_t$, so $\Delta p_t=(\phi+\theta)x_t-(\phi+\rho\theta)x_{t-1}+\dots$; $\rho$ from the
sign autocorrelation, then the two coefficients give $\theta$ and $\phi$. Moment conditions: the error orthogonal to $x_t$ and $x_{t-1}$.
\iqlookfor{The surprise in the order flow; identification through $\rho$.}
\end{solution}

\begin{solution}{iq:mx:decomposing-the-spread:4}
At one second the quotes have not absorbed the information: the one-second realised spread counts as profit what the next minute takes back. Look at
mark-outs at longer horizons, and at inventory losses and fees.
\iqlookfor{The horizon; mark-out curve; adverse selection arriving after the measurement.}
\end{solution}

\begin{solution}{iq:mx:decomposing-the-spread:5}
The permanent effect of a trade on the quote midpoint, as the cumulative impulse response of mid changes to a trade innovation, whatever the path. Lag
zero because quotes are revised after a trade within the interval; the ordering identifies the shock.
\iqlookfor{Permanent impact; impulse response; ordering assumption.}
\end{solution}

\begin{solution}{iq:mx:decomposing-the-spread:6}
Low-frequency estimators: Corwin--Schultz from highs and lows, Abdi--Ranaldo from closes and ranges, Roll from autocovariances. Check them on a later
period where intraday effective spreads are available, and on simulated data with known spreads; beware of liquid stocks where the spread is a small
share of the range.
\iqlookfor{A named estimator; a validation plan; its failure mode.}
\end{solution}
